% page 153 du fac-simile (image p-152 du scan)
% bloc 165.1 x 246.3 mm ; marge gauche 36.9 mm
\pgc{15.240mm}{0.000mm}{
\l{}
\l{}
\l{\cel{53}143.}
\l{}
\l{}
\l{\sou{eshafodo}. I. Konstrukturo qua subtenas la gradi por la spektanti.}
\l{\cel{3}- II. Konstrukturo sur placo publika, ube onu exekutas la}
\l{\cel{3}morto-verdikti. - III. Konstrukturo karpentura sur qua}
\l{\cel{3}laboristi okupesas da la konstrukto, la reparo o la dekoro}
\l{\cel{3}di la edifici. - DFIR.}
\l{}
\l{\sou{eshelono}. (\sou{armeo}).Personi, kozi, dispozita ye intervali. - F.}
\l{}
\l{\sou{-esk-}. Sufixo qua indikas \textquotedbl{}divenar\textquotedbl{}, \textquotedbl{}komencar...\textquotedbl{}.}
\l{}
\l{\sou{eskadro}. (\sou{Navaro}).Asemblajo de milito-navi komande da admiralo.}
\l{\cel{3}- DEFIRS.}
\l{}
\l{\sou{eskadrono}. (\sou{armeo)}. Trupo de armeani kavaliera ; dividuro de}
\l{\cel{3}regimento di la kavalrio, komandata da oficiro supra-grada.}
\l{\cel{3}- DEFIRS.}
\l{}
\l{\sou{eskalar}. (\sou{netrans.}) (\sou{pri navo}). Haltar en portuo, dum voyajo,}
\l{\cel{3}por embarkar o desembarkar pasajanti o vari. - FIS.}
\l{}
\l{\sou{eskaldar}. (\sou{trans.}) I. Brular per aquo varmega. - II. Imersar en}
\l{\cel{3}aquo bolianta. - EFS.}
\l{}
\l{\sou{eskalero}. (en\cel{1}\sou{edifico})\cel{2}Serio de gradi sequanta, fixigita, uzata}
\l{\cel{3}por acensar o decensar de etajo ad altra. - FIS.}
\l{}
\l{\sou{eskamotar}. (\sou{trans.}) Habile desaparigar objekto, per subtila pro-}
\l{\cel{3}cedo qua eskapas la regardo de la spektanto. - DFS.}
\l{}
\l{\sou{eskapar}. (\sou{trans.}) I. Tirar su de to, en quo onu retenesas. -}
\l{\cel{3}II. Cesar retenesar. - EFIS.}
\l{}
\l{\sou{eskaro.}\cel{1}(\sou{medic}.)\cel{2}Krusto nigratra, quan formacas la desorganizeso}
\l{\cel{3}di la tisui en parto atakata da gangreno o da la efiko di}
\l{\cel{3}korodivo. - EFIS.}
\l{}
\l{\sou{eskarbilo}. Fragmento de terkarbono, ne par-kombustita e mixita}
\l{\cel{3}a cindri. - FS.}
\l{}
\l{\sou{eskarolo}. (\sou{bot.}) Varietato di cikorio kultivata, quan onu manjas}
\l{\cel{3}kom salado. - L.\cel{1}\sou{lactuca scariola}. - DFIS.}
\l{}
\l{\sou{eskarpa}. Di qua la pento esas preske vertikala. - EFIS.}
\l{}
\l{\sou{eskartar}. (\sou{trans}.) I. Pozar (la parti di kozo) tale ke li kelke}
\l{\cel{3}interdistas. - II. Pozar kelka-diste de kozo, de persono.}
\l{\cel{3}- III. Forigar de la direciono quan onu devas sequar. - FIS.}
\l{}
\l{\sou{eskombro}. To quo restas de objekto ruptita. -FIS.}
\l{}
\l{eskopeto. Speco di funel-fusilo, quan onu +weris bandolieratre.}
\l{\cel{3}- FS.}
}
